\documentclass[12pt,a4paper]{report}

% --- Encodage et langue (XeLaTeX natif UTF-8) --------------------------------
\usepackage{fontspec}
\usepackage{polyglossia}
\setmainlanguage{french}

% --- Mise en page ------------------------------------------------------------
\usepackage[top=2.5cm, bottom=2.5cm, left=3cm, right=2.5cm]{geometry}
\usepackage{setspace}
\onehalfspacing

% --- Typographie -------------------------------------------------------------
\usepackage{microtype}
\usepackage{parskip}

% --- Couleurs et graphiques --------------------------------------------------
\usepackage[dvipsnames,table]{xcolor}
\usepackage{graphicx}
\usepackage{tikz}
\usetikzlibrary{shapes.geometric, arrows.meta, positioning, fit, backgrounds, calc}

% --- Tableaux ----------------------------------------------------------------
\usepackage{booktabs}
\usepackage{longtable}
\usepackage{tabularx}
\usepackage{multirow}
\usepackage{array}

% --- Code source -------------------------------------------------------------
\usepackage{listings}
\usepackage{lstautogobble}

% --- Liens -------------------------------------------------------------------
\usepackage[hidelinks]{hyperref}
\usepackage{amsmath}
\usepackage{cleveref}

% --- En-têtes ----------------------------------------------------------------
\usepackage{fancyhdr}
\usepackage{emptypage}

% --- Boîtes ------------------------------------------------------------------
\usepackage{tcolorbox}
\tcbuselibrary{skins, breakable, listings}

% --- Listes ------------------------------------------------------------------
\usepackage{enumitem}
\usepackage{titlesec}
\usepackage{pifont}

% ─────────────────────────────────────────────────────────────────────────────
% COULEURS
% ─────────────────────────────────────────────────────────────────────────────
\definecolor{fleetblue}{RGB}{19, 109, 236}
\definecolor{fleetdark}{RGB}{11, 17, 32}
\definecolor{fleetgray}{RGB}{107, 114, 128}
\definecolor{fleetsuccess}{RGB}{16, 185, 129}
\definecolor{fleetwarning}{RGB}{245, 158, 11}
\definecolor{fleeterror}{RGB}{239, 68, 68}
\definecolor{fleetpurple}{RGB}{139, 92, 246}
\definecolor{kernelcolor}{RGB}{6, 95, 70}
\definecolor{kernellight}{RGB}{209, 250, 229}
\definecolor{codebg}{RGB}{248, 249, 250}
\definecolor{codefg}{RGB}{36, 41, 47}
\definecolor{commentcolor}{RGB}{106, 153, 85}
\definecolor{keywordcolor}{RGB}{86, 156, 214}
\definecolor{stringcolor}{RGB}{206, 145, 120}

% ─────────────────────────────────────────────────────────────────────────────
% LISTINGS
% ─────────────────────────────────────────────────────────────────────────────
\lstdefinestyle{javaStyle}{
  language=Java,
  backgroundcolor=\color{codebg},
  basicstyle=\ttfamily\footnotesize\color{codefg},
  keywordstyle=\color{keywordcolor}\bfseries,
  stringstyle=\color{stringcolor},
  commentstyle=\color{commentcolor}\itshape,
  numberstyle=\tiny\color{fleetgray},
  numbers=left, numbersep=8pt,
  breaklines=true, frame=single,
  rulecolor=\color{fleetgray!40},
  xleftmargin=15pt, autogobble=true,
  morekeywords={UUID, Mono, Flux, record, var, WebClient}
}

\lstdefinestyle{yamlStyle}{
  language=bash,
  backgroundcolor=\color{codebg},
  basicstyle=\ttfamily\footnotesize\color{codefg},
  keywordstyle=\color{keywordcolor},
  stringstyle=\color{stringcolor},
  commentstyle=\color{commentcolor}\itshape,
  breaklines=true, frame=single,
  rulecolor=\color{fleetgray!40},
  xleftmargin=15pt, autogobble=true
}

% ─────────────────────────────────────────────────────────────────────────────
% BOITES
% ─────────────────────────────────────────────────────────────────────────────
\newtcolorbox{kernelbox}[2][]{
  colback=kernellight, colframe=kernelcolor,
  fonttitle=\bfseries\color{white},
  coltitle=white,
  attach boxed title to top left={yshift=-2mm, xshift=4mm},
  boxed title style={colback=kernelcolor},
  title={#2}, breakable, #1
}

\newtcolorbox{fleetbox}[2][]{
  colback=fleetblue!6, colframe=fleetblue,
  fonttitle=\bfseries\color{white},
  coltitle=white,
  attach boxed title to top left={yshift=-2mm, xshift=4mm},
  boxed title style={colback=fleetblue},
  title={#2}, breakable, #1
}

\newtcolorbox{warningbox}[1][]{
  colback=fleetwarning!8, colframe=fleetwarning,
  fonttitle=\bfseries\small,
  title={Attention}, breakable, #1
}

\newtcolorbox{infobox}[1][]{
  colback=fleetblue!5, colframe=fleetblue!60,
  fonttitle=\bfseries\small,
  title={Information}, breakable, #1
}

\newtcolorbox{contractbox}[1][]{
  colback=fleetdark!4, colframe=fleetdark,
  fonttitle=\bfseries\color{white},
  coltitle=white,
  title={Contrat d'intégration}, breakable, #1
}

% ─────────────────────────────────────────────────────────────────────────────
% TITRES
% ─────────────────────────────────────────────────────────────────────────────
\titleformat{\chapter}[display]
  {\normalfont\huge\bfseries\color{fleetdark}}
  {\chaptertitlename\ \thechapter}{20pt}{\Huge}
\titlespacing*{\chapter}{0pt}{-10pt}{30pt}

\titleformat{\section}
  {\normalfont\Large\bfseries\color{fleetblue}}
  {\thesection}{1em}{}
\titleformat{\subsection}
  {\normalfont\large\bfseries\color{fleetdark}}
  {\thesubsection}{1em}{}
\titleformat{\subsubsection}
  {\normalfont\normalsize\bfseries\color{fleetgray}}
  {\thesubsubsection}{1em}{}

% ─────────────────────────────────────────────────────────────────────────────
% EN-TETES
% ─────────────────────────────────────────────────────────────────────────────
\pagestyle{fancy}
\fancyhf{}
\fancyhead[L]{\small\textcolor{fleetgray}{\leftmark}}
\fancyhead[R]{\small\textcolor{fleetblue}{\textbf{FleetMan — Intégration Kernel RT-Comops}}}
\fancyfoot[C]{\small\textcolor{fleetgray}{\thepage}}
\renewcommand{\headrulewidth}{0.4pt}
\setlength{\headheight}{15pt}

% ─────────────────────────────────────────────────────────────────────────────
% DEBUT DU DOCUMENT
% ─────────────────────────────────────────────────────────────────────────────
\begin{document}

% ─────────────────────────────────────────────────────────────────────────────
% PAGE DE TITRE
% ─────────────────────────────────────────────────────────────────────────────
\begin{titlepage}
  \pagecolor{fleetdark}
  \color{white}
  \centering
  \vspace*{1.5cm}

  {\fontsize{50}{60}\selectfont\bfseries\textcolor{fleetblue}{FleetMan}}\\[0.3cm]
  {\Large\textcolor{white!70}{Plateforme de Gestion de Flotte Intelligente}}\\[2cm]

  \textcolor{fleetblue}{\rule{0.85\textwidth}{2pt}}\\[0.8cm]

  {\LARGE\bfseries Intégration du Kernel RT-Comops}\\[0.4cm]
  {\large Analyse, Identification des Modules Applicables et Plan d'Implémentation}\\[0.8cm]

  \textcolor{fleetblue}{\rule{0.85\textwidth}{1pt}}\\[1.5cm]

  \begin{tabular}{rl}
    \textcolor{white!60}{Projet :}       & \textbf{FleetMan — Module Fleet Management} \\[0.3cm]
    \textcolor{white!60}{Kernel :}       & \textbf{RT-Comops Backend — Architecture Multi-Tenant V3} \\[0.3cm]
    \textcolor{white!60}{Type :}         & \textbf{Document d'Intégration Technique} \\[0.3cm]
    \textcolor{white!60}{Modules Kernel :} & \textbf{20 cores analysés, 9 applicables à FleetMan} \\[0.3cm]
    \textcolor{white!60}{Version :}      & \textbf{1.0 — Juin 2026} \\[0.3cm]
    \textcolor{white!60}{Niveau :}       & \textbf{5\textsuperscript{ème} année Génie Informatique} \\
  \end{tabular}

  \vfill

  {\small\textcolor{white!50}{École Polytechnique — Département Informatique}}\\
  {\small\textcolor{white!50}{Document de Spécification d'Intégration}}
\end{titlepage}
\nopagecolor

\tableofcontents
\newpage

% ─────────────────────────────────────────────────────────────────────────────
\chapter*{Introduction Générale}
\addcontentsline{toc}{chapter}{Introduction Générale}
\thispagestyle{fancy}
% ─────────────────────────────────────────────────────────────────────────────

\section*{Contexte du Projet}

FleetMan est une plateforme de gestion de flotte de véhicules développée dans le cadre d'un projet de fin d'études de 5\textsuperscript{ème} année Génie Informatique. Elle constitue l'un des modules métier d'un écosystème applicatif plus vaste : la plateforme \textbf{RT-Comops}, un monolithe modulaire hexagonal multi-tenant développé par l'équipe de recherche et développement.

La plateforme RT-Comops est organisée autour d'un \textbf{Kernel} — le cœur technique transverse — qui concentre tous les mécanismes partagés entre les différents modules métier : authentification, gestion des identités, multi-tenancy, audit, événements, fichiers, organisations, rôles et permissions. Ce Kernel constitue la \textbf{base fondatrice} sur laquelle chaque module métier doit s'appuyer.

\section*{Objectif du Document}

Ce document a pour objectif de :

\begin{enumerate}[leftmargin=2cm]
  \item \textbf{Analyser} en profondeur les 20 cores du Kernel RT-Comops et comprendre leur rôle fonctionnel et technique.
  \item \textbf{Identifier} les modules du Kernel directement applicables et utiles au projet FleetMan.
  \item \textbf{Expliquer} comment chaque module Kernel sélectionné sera utilisé dans FleetMan, avec les patterns d'intégration concrets.
  \item \textbf{Planifier} l'implémentation de cette intégration de façon détaillée et progressive.
\end{enumerate}

\section*{Architecture Générale du Kernel}

Le Kernel RT-Comops est un monolithe modulaire hexagonal organisé en \textbf{20 cores fonctionnels}. Il repose sur la même stack technologique que FleetMan : Spring Boot WebFlux, Java 21, R2DBC, PostgreSQL, Redis, Kafka. Cette cohérence technologique facilite considérablement l'intégration.

\begin{kernelbox}{Les 20 Cores du Kernel RT-Comops}
\begin{enumerate}[leftmargin=1.5cm, itemsep=2pt]
  \item \textbf{kernel-core} — Multi-tenancy, JWT RS256, audit, outbox, ClientApplications
  \item \textbf{common-core} — Types partagés, BaseEntity, Address, Contact, PartyRef
  \item \textbf{actor-core} — Acteurs, BusinessActors, profils
  \item \textbf{auth-core} — Authentification, OIDC, OTP, MFA, JWT
  \item \textbf{roles-core} — Rôles, permissions, scopes RBAC
  \item \textbf{organization-core} — Organisations, agences, points d'intérêt
  \item \textbf{tp-core} — Tiers (clients, prospects, fournisseurs)
  \item \textbf{settings-core} — Paramètres métier, séquences documentaires
  \item \textbf{product-core} — Catalogue produits, variantes, prix
  \item \textbf{inventory-core} — Stocks, mouvements, inventaires
  \item \textbf{resource-core} — Ressources matérielles, véhicules, affectations
  \item \textbf{sales-core} — Commandes de vente
  \item \textbf{accounting-core} — Comptabilité générale
  \item \textbf{billing-core} — Facturation commerciale
  \item \textbf{treasury-core} — Trésorerie, comptes bancaires
  \item \textbf{cashier-core} — Caisse opérationnelle
  \item \textbf{administration-core} — Gouvernance, administration plateforme
  \item \textbf{file-core} — Gestion des fichiers et documents
  \item \textbf{hrm-core} — Ressources humaines
  \item \textbf{blockchain-core} — Wallets, signatures, ancrage documentaire
\end{enumerate}
\end{kernelbox}

% ─────────────────────────────────────────────────────────────────────────────
\chapter{Analyse Détaillée des 20 Cores du Kernel}
% ─────────────────────────────────────────────────────────────────────────────

\section{kernel-core — La Colonne Vertébrale Technique}

\subsection{Description et Rôle Principal}

Le \texttt{kernel-core} est le composant le plus fondamental de toute la plateforme. Il garantit l'exécution sûre de toutes les APIs en combinant :

\begin{itemize}[leftmargin=2cm]
  \item \textbf{Isolation multi-tenant logique} : chaque requête porte un \texttt{tenantId} injecté par le \texttt{TenantWebFilter} dans la chaîne réactive
  \item \textbf{Authentification des ClientApplications} : via les headers \texttt{X-Client-Id} et \texttt{X-Api-Key}
  \item \textbf{Validation JWT RS256} : extraction des claims \texttt{sub}, \texttt{tid}, \texttt{oid}, \texttt{aid}, \texttt{actor}
  \item \textbf{Outbox transactionnel} : écriture atomique des événements dans \texttt{kernel.outbox\_event}
  \item \textbf{Piste d'audit immuable} : \texttt{AuditTrail} en insertion seule
  \item \textbf{Quotas Redis} : limitation par \texttt{tenantId + clientId + serviceCode}
\end{itemize}

\subsection{Contrat d'Appel Minimal}

Tout backend qui se connecte au Kernel doit envoyer les headers suivants :

\begin{lstlisting}[style=yamlStyle, caption={Headers obligatoires vers le Kernel}]
X-Client-Id: <identifiant stable de la ClientApplication>
X-Api-Key: <secret serveur de l'application>
X-Tenant-Id: <tenant cible de l'appel>
Authorization: Bearer <accessToken>  # requis pour actions utilisateur
X-Organization-Id: <requis pour services metier scopes organisation>
\end{lstlisting}

\subsection{Catalogue des Services Autorisés}

Le Kernel contrôle l'accès aux services via le catalogue \texttt{PlatformServiceCode} :

\begin{longtable}{p{4cm}p{9cm}}
\toprule
\textbf{Code Service} & \textbf{Description} \\
\midrule
\endhead
\texttt{ORGANIZATION} & Gestion des organisations et agences \\
\texttt{SETTINGS} & Paramètres métier et séquences \\
\texttt{COMMERCIAL} & Tiers, prospects, clients \\
\texttt{PRODUCT} & Catalogue produits \\
\texttt{INVENTORY} & Gestion des stocks \\
\texttt{SALES} & Commandes de vente \\
\texttt{BILLING} & Facturation commerciale \\
\texttt{ACCOUNTING} & Comptabilité générale \\
\texttt{BANKING} & Comptes bancaires \\
\texttt{TREASURY} & Trésorerie \\
\texttt{CASHIER} & Caisse opérationnelle \\
\texttt{RESOURCE} & Ressources matérielles \\
\texttt{HRM} & Ressources humaines \\
\texttt{BLOCKCHAIN} & Wallets et ancrage \\
\bottomrule
\caption{Catalogue PlatformServiceCode du Kernel}
\end{longtable}

\section{common-core — Socle de Types Partagés}

\subsection{Description}

Le \texttt{common-core} fournit les abstractions réutilisées par tous les autres cores :

\begin{itemize}[leftmargin=2cm]
  \item \textbf{BaseEntity} : classe abstraite avec \texttt{UUID id}, \texttt{tenantId}, \texttt{createdAt}, \texttt{updatedAt}
  \item \textbf{Address} : adresse polymorphique attachable à toute entité (lignes, ville, pays, GPS)
  \item \textbf{Contact} : contact polymorphique avec vérification email/téléphone
  \item \textbf{PartyRef} : référence légère inter-cores (\texttt{partyType} + \texttt{partyId})
  \item \textbf{ApiResponse<T>} et \textbf{PageResult<T>} : enveloppes de réponse JSON uniformes
  \item \textbf{PlatformServiceCode} : énumération partagée des services
\end{itemize}

\subsection{Invariants Clés}

\begin{itemize}[leftmargin=2cm]
  \item Une seule adresse par défaut par combinaison (\texttt{addressableId}, \texttt{addressableType}, \texttt{addressType})
  \item \texttt{PartyRef.partyType} limité à \texttt{ACTOR} et \texttt{ORGANIZATION}
  \item \texttt{tenantId} obligatoire sur toutes les structures polymorphiques
\end{itemize}

\section{actor-core — Gestion des Acteurs}

\subsection{Description}

L'\texttt{actor-core} gère les acteurs de la plateforme : personnes physiques (\texttt{Actor}), entités commerciales (\texttt{BusinessActor}), et leurs profils détaillés. Il couvre :

\begin{itemize}[leftmargin=2cm]
  \item Création, mise à jour et suppression douce des acteurs
  \item Rattachement d'un acteur à une organisation
  \item Cycle de vie : \texttt{PENDING} → \texttt{ACTIVE} → \texttt{SUSPENDED} → \texttt{DELETED}
  \item Réactivation par auto-service ou approbation administrative
\end{itemize}

\section{auth-core — Authentification et Identité}

\subsection{Description}

L'\texttt{auth-core} gère l'ensemble du cycle d'authentification :

\begin{itemize}[leftmargin=2cm]
  \item Création de comptes utilisateurs par un administrateur IAM
  \item Authentification locale tenant-scoped avec émission de JWT RS256
  \item Découverte globale des contextes de login (multi-tenant)
  \item Sélection explicite d'un contexte tenant/organisation
  \item Identification email-first et inscription contextuelle
  \item Réinitialisation de mot de passe (SMTP ou PREVIEW\_ONLY)
  \item Vérification d'email post-authentification
  \item Gestion du profil utilisateur (\texttt{/users/me})
\end{itemize}

\subsection{Flow JWT RS256}

Le Kernel émet des JWT RS256 signés avec une clé privée. Les backends consommateurs vérifient la signature via le document \texttt{/.well-known/jwks.json}. Les claims transportés sont : \texttt{sub} (userId), \texttt{tid} (tenantId), \texttt{oid} (organizationId), \texttt{aid} (agencyId), \texttt{actor} (actorId).

\section{roles-core — Rôles et Permissions RBAC}

\subsection{Description}

Le \texttt{roles-core} implémente un système RBAC (Role-Based Access Control) complet :

\begin{itemize}[leftmargin=2cm]
  \item Création de rôles avec codes de permission granulaires
  \item Affectation de rôles à des utilisateurs avec portée (TENANT, ORGANIZATION, AGENCY)
  \item Résolution des permissions avec cache Redis optionnel
  \item Construction des autorités scopées à partir d'une affectation
  \item Invalidation du cache après modification d'affectation
  \item Provisionnement de templates de rôles (réservé à \texttt{administration-core})
\end{itemize}

\section{organization-core — Organisations et Agences}

\subsection{Description}

L'\texttt{organization-core} gère la structure organisationnelle :

\begin{itemize}[leftmargin=2cm]
  \item Création d'une organisation avec agence siège automatique
  \item Création d'agences opérationnelles
  \item Gestion des horaires d'ouverture
  \item Lien fonctionnel acteur-organisation
  \item Suspension d'une organisation et de ses agences
  \item Ajout de points d'intérêt géographiques (GPS)
\end{itemize}

\section{tp-core — Gestion des Tiers}

\subsection{Description}

Le \texttt{tp-core} gère les tiers commerciaux (clients, prospects, fournisseurs) :

\begin{itemize}[leftmargin=2cm]
  \item Création d'un tiers avec rôle client
  \item Qualification commerciale d'un prospect
  \item Conversion prospect → client
  \item Enregistrement d'interactions commerciales
  \item Ajout de comptes bancaires
  \item Recherche plein-texte (Elasticsearch / fallback PostgreSQL)
\end{itemize}

\section{settings-core — Paramètres Métier}

\subsection{Description}

Le \texttt{settings-core} gère la configuration métier de la plateforme :

\begin{itemize}[leftmargin=2cm]
  \item Génération de numéros documentaires avec fallback agence → organisation
  \item Lecture des paramètres métier effectifs
  \item Mise à jour des options générales d'une organisation
  \item Gestion des séquences documentaires par agence
  \item Politique opérationnelle d'agence avec repli organisation
\end{itemize}

\section{resource-core — Ressources Matérielles}

\subsection{Description}

Le \texttt{resource-core} est le module le plus directement lié à FleetMan. Il gère les ressources matérielles de l'organisation :

\begin{itemize}[leftmargin=2cm]
  \item Enregistrement d'une ressource matérielle (véhicule, équipement)
  \item Affectation d'une ressource à un acteur (conducteur)
  \item Envoi en maintenance et retour à l'état opérationnel
  \item Mise à jour de la localisation GPS en temps réel
  \item Mise au rebut irréversible
\end{itemize}

\section{file-core — Gestion des Fichiers}

\subsection{Description}

Le \texttt{file-core} gère le cycle de vie des fichiers et documents :

\begin{itemize}[leftmargin=2cm]
  \item Téléversement avec audit système
  \item Téléchargement sécurisé
  \item Rattachement d'un document à une cible métier (polymorphique)
  \item Consultation des documents d'une cible
  \item Politique de gouvernance documentaire (expiration, révision)
\end{itemize}

\section{accounting-core, billing-core, treasury-core — Finance}

Ces trois cores couvrent le cycle financier complet :

\begin{itemize}[leftmargin=2cm]
  \item \textbf{accounting-core} : Comptabilité générale, factures canoniques, règlements, postings
  \item \textbf{billing-core} : Facturation commerciale legacy, bons de livraison, paiements
  \item \textbf{treasury-core} : Comptes bancaires, relevés, rapprochements, chèques
\end{itemize}

\section{hrm-core — Ressources Humaines}

Le \texttt{hrm-core} couvre l'ensemble du cycle RH : employés, contrats, congés, paie, recrutement, onboarding, formation, frais, missions, médical et KPI RH. Il s'intègre avec \texttt{actor-core} et \texttt{organization-core}.

\section{blockchain-core — Ancrage et Wallets}

Le \texttt{blockchain-core} fournit une blockchain interne pour : wallets, signatures numériques, blocs, ancrage documentaire et ancrage d'événements métier. Il garantit l'immuabilité et la traçabilité des opérations critiques.

\section{administration-core — Gouvernance}

L'\texttt{administration-core} gère la gouvernance de la plateforme : approbation des BusinessActors, mise à jour des options de plateforme avec audit, consultation de la piste d'audit, supervision de l'état du système.

% ─────────────────────────────────────────────────────────────────────────────
\chapter{Identification des Modules Kernel Applicables à FleetMan}
% ─────────────────────────────────────────────────────────────────────────────

\section{Méthodologie de Sélection}

Pour identifier les modules Kernel applicables à FleetMan, nous avons appliqué trois critères de sélection :

\begin{enumerate}[leftmargin=2cm]
  \item \textbf{Pertinence fonctionnelle} : le module couvre un besoin déjà présent ou planifié dans FleetMan
  \item \textbf{Évitement de la duplication} : FleetMan ne doit pas réimplémenter ce que le Kernel fournit déjà
  \item \textbf{Cohérence architecturale} : l'intégration respecte le contrat backend-to-kernel
\end{enumerate}

\section{Tableau de Sélection des Modules}

\begin{longtable}{p{3.5cm}p{2cm}p{7.5cm}}
\toprule
\textbf{Module Kernel} & \textbf{Applicable} & \textbf{Justification} \\
\midrule
\endhead
\texttt{kernel-core} & \textcolor{fleetsuccess}{\textbf{OUI}} & Fondation obligatoire : multi-tenancy, JWT, audit, outbox \\
\texttt{common-core} & \textcolor{fleetsuccess}{\textbf{OUI}} & BaseEntity, Address, Contact, PageResult réutilisables \\
\texttt{actor-core} & \textcolor{fleetsuccess}{\textbf{OUI}} & Gestionnaires et conducteurs sont des acteurs \\
\texttt{auth-core} & \textcolor{fleetsuccess}{\textbf{OUI}} & Remplace l'authentification locale actuelle de FleetMan \\
\texttt{roles-core} & \textcolor{fleetsuccess}{\textbf{OUI}} & Remplace le RBAC local (FLEET\_MANAGER, FLEET\_DRIVER) \\
\texttt{organization-core} & \textcolor{fleetsuccess}{\textbf{OUI}} & Une flotte = une organisation, un dépôt = une agence \\
\texttt{tp-core} & \textcolor{fleetsuccess}{\textbf{OUI}} & Les clients des missions sont des tiers \\
\texttt{settings-core} & \textcolor{fleetsuccess}{\textbf{OUI}} & Numéros de missions, paramètres métier \\
\texttt{resource-core} & \textcolor{fleetsuccess}{\textbf{OUI}} & Les véhicules sont des ressources matérielles \\
\texttt{file-core} & \textcolor{fleetsuccess}{\textbf{OUI}} & Documents légaux véhicules/conducteurs \\
\texttt{billing-core} & \textcolor{fleetwarning}{\textbf{PARTIEL}} & Facturation des missions (Module 7 du cahier) \\
\texttt{accounting-core} & \textcolor{fleetwarning}{\textbf{PARTIEL}} & Comptabilisation des dépenses opérationnelles \\
\texttt{product-core} & \textcolor{fleeterror}{\textbf{NON}} & Pas de catalogue produits dans FleetMan \\
\texttt{inventory-core} & \textcolor{fleeterror}{\textbf{NON}} & Pas de gestion de stocks dans FleetMan \\
\texttt{sales-core} & \textcolor{fleeterror}{\textbf{NON}} & Pas de commandes de vente au sens strict \\
\texttt{treasury-core} & \textcolor{fleeterror}{\textbf{NON}} & Hors périmètre FleetMan \\
\texttt{cashier-core} & \textcolor{fleeterror}{\textbf{NON}} & Hors périmètre FleetMan \\
\texttt{administration-core} & \textcolor{fleetwarning}{\textbf{PARTIEL}} & Gouvernance du super-admin FleetMan \\
\texttt{hrm-core} & \textcolor{fleeterror}{\textbf{NON}} & Hors périmètre FleetMan (RH complète) \\
\texttt{blockchain-core} & \textcolor{fleetwarning}{\textbf{FUTUR}} & Ancrage des incidents critiques (phase future) \\
\bottomrule
\caption{Sélection des modules Kernel applicables à FleetMan}
\end{longtable}

\section{Les 9 Modules Retenus}

\begin{tcolorbox}[colback=kernellight, colframe=kernelcolor, title={\bfseries\color{white} Modules Kernel retenus pour FleetMan}]
\begin{enumerate}[leftmargin=1.5cm]
  \item \textbf{kernel-core} — Fondation obligatoire (multi-tenancy, JWT, audit, outbox)
  \item \textbf{common-core} — Types partagés (BaseEntity, Address, Contact, PageResult)
  \item \textbf{actor-core} — Gestion des acteurs (gestionnaires, conducteurs)
  \item \textbf{auth-core} — Authentification centralisée (remplace l'auth locale)
  \item \textbf{roles-core} — RBAC centralisé (remplace les rôles locaux)
  \item \textbf{organization-core} — Structure organisationnelle (flotte = organisation)
  \item \textbf{tp-core} — Gestion des tiers (clients des missions)
  \item \textbf{settings-core} — Paramètres métier et numérotation
  \item \textbf{resource-core} — Ressources matérielles (véhicules)
  \item \textbf{file-core} — Documents légaux (assurance, permis, carte grise)
\end{enumerate}
\end{tcolorbox}

% ─────────────────────────────────────────────────────────────────────────────
\chapter{Analyse Détaillée de Chaque Intégration}
% ─────────────────────────────────────────────────────────────────────────────

\section{Intégration 1 — kernel-core : Fondation Multi-Tenant}

\subsection{Situation Actuelle dans FleetMan}

Actuellement, FleetMan gère son propre contexte de sécurité via \texttt{SecurityConfig}, \texttt{JwtAuthenticationManager} et \texttt{BearerTokenServerAuthenticationConverter}. L'authentification est déléguée à un service Auth externe via \texttt{RemoteAuthAdapter}, mais sans le contrat formel du Kernel.

\subsection{Ce que le kernel-core Apporte}

\begin{itemize}[leftmargin=2cm]
  \item \textbf{Multi-tenancy logique} : FleetMan devient un tenant du Kernel. Chaque entreprise de transport est un tenant distinct, avec isolation complète des données.
  \item \textbf{ClientApplication} : FleetMan s'enregistre comme \texttt{fleet-management-backend} avec ses \texttt{allowedServiceCodes} : \texttt{ORGANIZATION}, \texttt{RESOURCE}, \texttt{COMMERCIAL}, \texttt{SETTINGS}, \texttt{FILE}.
  \item \textbf{JWT RS256 centralisé} : FleetMan vérifie les tokens via le JWKS du Kernel au lieu de son propre mécanisme.
  \item \textbf{Audit immuable} : toutes les opérations critiques (création véhicule, incident, maintenance) sont tracées dans \texttt{kernel.outbox\_event}.
  \item \textbf{Outbox transactionnel} : les événements FleetMan (incident critique, maintenance créée) sont publiés via l'outbox du Kernel vers Kafka.
\end{itemize}

\subsection{Mapping Conceptuel}

\begin{longtable}{p{5cm}p{5cm}p{3cm}}
\toprule
\textbf{Concept FleetMan} & \textbf{Concept Kernel} & \textbf{Mapping} \\
\midrule
\endhead
Entreprise de transport & Tenant & 1:1 \\
Application FleetMan & ClientApplication & 1:1 \\
Token JWT actuel & JWT RS256 Kernel & Remplacement \\
Audit local & AuditTrail Kernel & Remplacement \\
Événements Kafka locaux & Outbox Kernel & Enrichissement \\
\bottomrule
\caption{Mapping kernel-core vers FleetMan}
\end{longtable}

\subsection{Implémentation Technique}

\begin{lstlisting}[style=javaStyle, caption={Configuration de la ClientApplication FleetMan}]
// Dans application.yml
application:
  kernel:
    url: https://kernel.rt-comops.com
    client-id: fleet-management-backend
    api-key: ${FLEET_KERNEL_API_KEY}
    tenant-id: ${FLEET_TENANT_ID}
    allowed-services:
      - ORGANIZATION
      - RESOURCE
      - COMMERCIAL
      - SETTINGS
      - FILE
\end{lstlisting}

\begin{lstlisting}[style=javaStyle, caption={Adapter WebClient vers le Kernel}]
@Component
@RequiredArgsConstructor
public class KernelWebClientAdapter {

    private final WebClient kernelWebClient;

    @Value("${application.kernel.client-id}")
    private String clientId;

    @Value("${application.kernel.api-key}")
    private String apiKey;

    @Value("${application.kernel.tenant-id}")
    private String tenantId;

    public WebClient.RequestHeadersSpec<?> withKernelHeaders(
            WebClient.RequestHeadersSpec<?> spec,
            String organizationId,
            String bearerToken) {
        return spec
            .header("X-Client-Id", clientId)
            .header("X-Api-Key", apiKey)
            .header("X-Tenant-Id", tenantId)
            .header("X-Organization-Id", organizationId)
            .header("Authorization", "Bearer " + bearerToken);
    }
}
\end{lstlisting}

\section{Intégration 2 — common-core : Types Partagés}

\subsection{Situation Actuelle}

FleetMan utilise ses propres records Java immuables (\texttt{Vehicle}, \texttt{Driver}, \texttt{Fleet}) sans base commune. La pagination est gérée par le \texttt{PageResponse<T>} créé dans le Chapitre 1 de notre implémentation.

\subsection{Ce que le common-core Apporte}

\begin{itemize}[leftmargin=2cm]
  \item \textbf{BaseEntity} : FleetMan adopte la classe abstraite \texttt{BaseEntity} du Kernel comme base de toutes ses entités, garantissant la présence de \texttt{tenantId}, \texttt{createdAt}, \texttt{updatedAt}.
  \item \textbf{Address polymorphique} : les adresses des flottes, dépôts et conducteurs utilisent le modèle \texttt{Address} du Kernel au lieu de champs texte dispersés.
  \item \textbf{Contact polymorphique} : les contacts des gestionnaires et conducteurs utilisent \texttt{Contact} du Kernel.
  \item \textbf{PageResult<T>} : notre \texttt{PageResponse<T>} est aligné sur le format \texttt{PageResult<T>} du Kernel pour une cohérence API globale.
  \item \textbf{ApiResponse<T>} : toutes les réponses FleetMan adoptent l'enveloppe standard du Kernel.
\end{itemize}

\subsection{Mapping Conceptuel}

\begin{longtable}{p{5cm}p{5cm}p{3cm}}
\toprule
\textbf{Concept FleetMan} & \textbf{Concept Kernel} & \textbf{Action} \\
\midrule
\endhead
Champs \texttt{createdAt/updatedAt} & \texttt{BaseEntity} & Héritage \\
Adresse de la flotte & \texttt{Address} polymorphique & Remplacement \\
Contact du gestionnaire & \texttt{Contact} polymorphique & Remplacement \\
\texttt{PageResponse<T>} & \texttt{PageResult<T>} & Alignement \\
Réponses REST & \texttt{ApiResponse<T>} & Adoption \\
\bottomrule
\caption{Mapping common-core vers FleetMan}
\end{longtable}

\section{Intégration 3 — auth-core : Authentification Centralisée}

\subsection{Situation Actuelle}

FleetMan délègue déjà l'authentification à un service Auth externe via \texttt{RemoteAuthAdapter} et \texttt{AuthApiClient}. Le \texttt{JwtAuthenticationManager} valide les tokens en interrogeant ce service.

\subsection{Ce que l'auth-core Apporte}

L'\texttt{auth-core} du Kernel \textbf{remplace complètement} le service Auth externe actuel. Les bénéfices sont :

\begin{itemize}[leftmargin=2cm]
  \item \textbf{Login multi-contexte} : un utilisateur peut avoir des comptes dans plusieurs tenants (plusieurs entreprises de transport)
  \item \textbf{Découverte des contextes} : l'endpoint \texttt{/auth/contexts} liste les tenants disponibles pour un email donné
  \item \textbf{JWT RS256 standardisé} : les claims \texttt{tid}, \texttt{oid}, \texttt{aid} sont automatiquement présents
  \item \textbf{OTP et MFA} : sécurité renforcée pour les gestionnaires de flotte
  \item \textbf{Réinitialisation de mot de passe} : flux complet avec SMTP
  \item \textbf{Vérification d'email} : obligatoire pour les nouveaux comptes
\end{itemize}

\subsection{Mapping des Rôles}

\begin{longtable}{p{5cm}p{5cm}p{3cm}}
\toprule
\textbf{Rôle FleetMan Actuel} & \textbf{Rôle Kernel} & \textbf{Scope} \\
\midrule
\endhead
\texttt{SUPER\_ADMIN} & \texttt{SYSTEM\_ADMIN} & SYSTEM \\
\texttt{FLEET\_ADMIN} & \texttt{TENANT\_ADMIN} & TENANT \\
\texttt{FLEET\_MANAGER} & \texttt{ORG\_MANAGER} & ORGANIZATION \\
\texttt{FLEET\_DRIVER} & \texttt{ORG\_OPERATOR} & ORGANIZATION \\
\bottomrule
\caption{Mapping des rôles FleetMan vers le Kernel}
\end{longtable}

\subsection{Adaptation du SecurityConfig}

\begin{lstlisting}[style=javaStyle, caption={Nouveau SecurityConfig avec JWT Kernel RS256}]
@Configuration
@EnableWebFluxSecurity
@EnableReactiveMethodSecurity
public class SecurityConfig {

    @Bean
    public SecurityWebFilterChain filterChain(ServerHttpSecurity http,
            KernelJwtAuthenticationManager authManager) {
        return http
            .cors(cors -> cors.configurationSource(corsSource()))
            .csrf(ServerHttpSecurity.CsrfSpec::disable)
            .authorizeExchange(ex -> ex
                .pathMatchers("/v3/api-docs/**", "/swagger-ui/**",
                              "/actuator/**", "/api/v1/auth/**").permitAll()
                .anyExchange().authenticated()
            )
            // Validation JWT via JWKS du Kernel
            .oauth2ResourceServer(oauth2 -> oauth2
                .jwt(jwt -> jwt.authenticationManager(authManager))
            )
            .build();
    }
}
\end{lstlisting}

\section{Intégration 4 — roles-core : RBAC Centralisé}

\subsection{Situation Actuelle}

FleetMan utilise \texttt{@PreAuthorize("hasRole('FLEET\_MANAGER')")} avec des rôles définis localement dans le JWT. Le \texttt{JwtAuthenticationManager} extrait les rôles du token.

\subsection{Ce que le roles-core Apporte}

\begin{itemize}[leftmargin=2cm]
  \item \textbf{Rôles avec portée} : un rôle peut être limité à une organisation ou une agence spécifique
  \item \textbf{Permissions granulaires} : au lieu de rôles grossiers, des codes de permission précis (\texttt{VEHICLE\_CREATE}, \texttt{DRIVER\_MANAGE}, \texttt{INCIDENT\_REPORT})
  \item \textbf{Cache Redis} : les permissions sont mises en cache pour des performances optimales
  \item \textbf{Templates de rôles} : des rôles prédéfinis pour les cas d'usage FleetMan
\end{itemize}

\subsection{Définition des Rôles FleetMan dans le Kernel}

\begin{lstlisting}[style=yamlStyle, caption={Templates de rôles FleetMan à provisionner dans le Kernel}]
roles:
  - code: FLEET_MANAGER
    name: "Gestionnaire de Flotte"
    permissions:
      - VEHICLE_CREATE
      - VEHICLE_READ
      - VEHICLE_UPDATE
      - VEHICLE_DELETE
      - DRIVER_CREATE
      - DRIVER_READ
      - DRIVER_MANAGE
      - FLEET_MANAGE
      - INCIDENT_READ
      - MAINTENANCE_MANAGE
      - ASSIGNMENT_MANAGE
      - DOCUMENT_MANAGE
      - KPI_READ
      - EXPENSE_APPROVE

  - code: FLEET_DRIVER
    name: "Conducteur"
    permissions:
      - VEHICLE_READ
      - TRIP_START
      - TRIP_END
      - TELEMETRY_SEND
      - INCIDENT_REPORT
      - MAINTENANCE_DECLARE
      - FUEL_RECHARGE_CREATE
      - ASSIGNMENT_VIEW
      - DOCUMENT_VIEW
\end{lstlisting}

\section{Intégration 5 — organization-core : Flotte = Organisation}

\subsection{Mapping Conceptuel Clé}

C'est l'une des intégrations les plus structurantes. Le mapping est le suivant :

\begin{longtable}{p{5cm}p{5cm}p{3cm}}
\toprule
\textbf{Concept FleetMan} & \textbf{Concept Kernel} & \textbf{Cardinalité} \\
\midrule
\endhead
Entreprise de transport & Tenant & 1:1 \\
Flotte de véhicules & Organisation & 1:1 \\
Dépôt / Garage & Agence & N:1 \\
Gestionnaire de flotte & Acteur lié à l'organisation & N:1 \\
Zone géographique & Point d'intérêt d'agence & N:1 \\
\bottomrule
\caption{Mapping organization-core vers FleetMan}
\end{longtable}

\subsection{Conséquences Architecturales}

\begin{itemize}[leftmargin=2cm]
  \item La table \texttt{fleet.fleets} devient une \textbf{projection locale} de l'organisation Kernel, enrichie des données spécifiques à FleetMan (nombre de véhicules, etc.)
  \item L'\texttt{organizationId} du Kernel est propagé dans tous les appels FleetMan via le header \texttt{X-Organization-Id}
  \item La création d'une flotte dans FleetMan déclenche la création d'une organisation dans le Kernel (via l'outbox)
  \item La suspension d'une organisation dans le Kernel suspend automatiquement la flotte correspondante
\end{itemize}

\subsection{Synchronisation via Outbox}

\begin{lstlisting}[style=javaStyle, caption={Synchronisation flotte-organisation via événements Kafka}]
// Consommateur Kafka dans FleetMan
@Component
@RequiredArgsConstructor
@Slf4j
public class OrganizationEventConsumer {

    private final FleetRepositoryPort fleetPort;

    @KafkaListener(topics = "iwm.events.business",
                   groupId = "fleet-management-group")
    public void onOrganizationEvent(OrganizationEvent event) {
        if ("ORGANIZATION_SUSPENDED".equals(event.type())) {
            // Suspendre la flotte correspondante
            fleetPort.findByOrganizationId(event.organizationId())
                .flatMap(fleet -> fleetPort.suspend(fleet.id()))
                .subscribe();
        }
    }
}
\end{lstlisting}

\section{Intégration 6 — tp-core : Clients des Missions}

\subsection{Situation Actuelle}

Le Module 7 (Clients et Missions) de notre cahier d'analyse prévoit une table \texttt{fleet.clients} locale. Avec le Kernel, ces clients deviennent des \textbf{tiers} gérés par \texttt{tp-core}.

\subsection{Ce que le tp-core Apporte}

\begin{itemize}[leftmargin=2cm]
  \item \textbf{Cycle de vie complet} : prospect → client avec qualification commerciale
  \item \textbf{Recherche plein-texte} : via Elasticsearch pour retrouver rapidement un client
  \item \textbf{Comptes bancaires} : pour la facturation des missions
  \item \textbf{Historique des interactions} : toutes les missions sont des interactions commerciales
  \item \textbf{Déduplication} : un client peut être partagé entre plusieurs flottes du même tenant
\end{itemize}

\subsection{Adaptation du Module 7}

La table \texttt{fleet.clients} locale est remplacée par une \textbf{référence légère} vers le tiers Kernel :

\begin{lstlisting}[style=javaStyle, caption={Adaptation de la table missions avec référence tiers Kernel}]
-- Adaptation de fleet.missions
ALTER TABLE fleet.missions
    ADD COLUMN IF NOT EXISTS third_party_id UUID,  -- ID du tiers dans le Kernel
    ADD COLUMN IF NOT EXISTS third_party_name VARCHAR(255); -- Dénormalisé
-- La colonne client_id locale devient optionnelle (migration progressive)
\end{lstlisting}

\section{Intégration 7 — settings-core : Paramètres et Numérotation}

\subsection{Ce que le settings-core Apporte à FleetMan}

\begin{itemize}[leftmargin=2cm]
  \item \textbf{Numéros de missions} : génération automatique de références uniques (\texttt{MSN-2026-001234}) avec fallback agence → organisation
  \item \textbf{Numéros de factures} : séquences documentaires pour les factures de missions
  \item \textbf{Paramètres métier} : taux de TVA, devise par défaut (XAF), limites opérationnelles
  \item \textbf{Politique opérationnelle} : horaires de service, zones autorisées par agence
\end{itemize}

\subsection{Utilisation dans FleetMan}

\begin{lstlisting}[style=javaStyle, caption={Appel au settings-core pour générer un numéro de mission}]
@Component
@RequiredArgsConstructor
public class MissionNumberGenerator {

    private final KernelWebClientAdapter kernelClient;

    public Mono<String> generateMissionNumber(String organizationId,
                                               String agencyId,
                                               String bearerToken) {
        return kernelClient.withKernelHeaders(
            kernelWebClient.post()
                .uri("/api/v1/settings/sequences/next")
                .bodyValue(Map.of(
                    "sequenceCode", "MISSION",
                    "agencyId", agencyId
                )),
            organizationId, bearerToken
        ).retrieve()
         .bodyToMono(String.class);
    }
}
\end{lstlisting}

\section{Intégration 8 — resource-core : Véhicules comme Ressources}

\subsection{Situation Actuelle}

FleetMan gère ses véhicules localement via \texttt{VehicleService} et \texttt{VehiclePersistenceAdapter}. Il existe déjà un \texttt{ExternalVehiclePort} et un \texttt{VehicleApiAdapter} qui communiquent avec un service véhicule externe.

\subsection{Ce que le resource-core Apporte}

Le \texttt{resource-core} du Kernel \textbf{remplace} le service véhicule externe actuel. Les véhicules deviennent des \textbf{ressources matérielles} du Kernel :

\begin{itemize}[leftmargin=2cm]
  \item \textbf{Enregistrement centralisé} : un véhicule est enregistré une fois dans le Kernel et référencé dans FleetMan
  \item \textbf{Affectation standardisée} : l'affectation conducteur-véhicule utilise le mécanisme d'affectation acteur-ressource du Kernel
  \item \textbf{Cycle de vie unifié} : AVAILABLE → ON\_TRIP → MAINTENANCE → DECOMMISSIONED
  \item \textbf{Localisation GPS} : mise à jour de la position via le port dédié du Kernel
  \item \textbf{Maintenance intégrée} : le passage en maintenance est synchronisé avec le Kernel
\end{itemize}

\subsection{Mapping Technique}

\begin{longtable}{p{5cm}p{5cm}p{3cm}}
\toprule
\textbf{Concept FleetMan} & \textbf{Concept Kernel} & \textbf{Action} \\
\midrule
\endhead
\texttt{Vehicle} & \texttt{Resource} (type VEHICLE) & Synchronisation \\
\texttt{VehicleService.create()} & \texttt{resource-core.register()} & Délégation \\
\texttt{VehicleService.assignDriver()} & \texttt{resource-core.assign()} & Délégation \\
\texttt{TripService.startTrip()} & \texttt{resource-core.updateGps()} & Enrichissement \\
\texttt{MaintenanceService.create()} & \texttt{resource-core.sendToMaintenance()} & Synchronisation \\
\bottomrule
\caption{Mapping resource-core vers FleetMan}
\end{longtable}

\subsection{Adaptation du VehicleService}

\begin{lstlisting}[style=javaStyle, caption={VehicleService enrichi avec resource-core}]
@Service
@RequiredArgsConstructor
@Slf4j
public class VehicleService implements ManageVehicleUseCase {

    private final VehiclePersistencePort vehiclePort;
    private final KernelResourceAdapter kernelResourceAdapter;

    @Override
    public Mono<Vehicle> createIndependentVehicle(VehicleRequest request,
                                                   UUID managerId,
                                                   String token) {
        // 1. Enregistrer dans le Kernel resource-core
        return kernelResourceAdapter.registerResource(
                ResourceRegistrationRequest.builder()
                    .resourceType("VEHICLE")
                    .name(request.brand() + " " + request.model())
                    .serialNumber(request.vehicleSerialNumber())
                    .organizationId(managerId.toString())
                    .build(), token
        )
        // 2. Sauvegarder localement avec l'ID Kernel
        .flatMap(kernelResource -> {
            Vehicle vehicle = buildVehicle(request, managerId,
                                           kernelResource.id());
            return vehiclePort.saveLocalData(vehicle);
        });
    }
}
\end{lstlisting}

\section{Intégration 9 — file-core : Documents Légaux}

\subsection{Situation Actuelle}

Le Module 2 (Documents Légaux) de notre implémentation stocke les URLs des fichiers dans les tables \texttt{fleet.vehicle\_documents} et \texttt{fleet.driver\_documents}. Les fichiers eux-mêmes sont stockés localement ou sur un service externe non standardisé.

\subsection{Ce que le file-core Apporte}

\begin{itemize}[leftmargin=2cm]
  \item \textbf{Stockage centralisé} : tous les fichiers (assurance, permis, carte grise) sont gérés par le Kernel
  \item \textbf{Rattachement polymorphique} : un document est attaché à une cible métier (véhicule ou conducteur) via \texttt{targetType} + \texttt{targetId}
  \item \textbf{Politique de gouvernance} : expiration automatique, révision périodique
  \item \textbf{Audit de téléversement} : chaque upload est tracé dans l'audit Kernel
  \item \textbf{Téléchargement sécurisé} : accès contrôlé par les permissions RBAC
\end{itemize}

\subsection{Adaptation du DocumentService}

\begin{lstlisting}[style=javaStyle, caption={DocumentService enrichi avec file-core}]
@Override
public Mono<VehicleDocument> addVehicleDocument(
        AddVehicleDocumentCommand cmd) {

    // 1. Téléverser le fichier dans le Kernel file-core
    return kernelFileAdapter.uploadFile(
            cmd.fileContent(), cmd.fileName(),
            "VEHICLE_DOCUMENT", cmd.vehicleId().toString()
    )
    // 2. Récupérer l'URL sécurisée du Kernel
    .flatMap(kernelFile -> {
        VehicleDocument doc = new VehicleDocument(
            null, cmd.vehicleId(),
            VehicleDocument.DocType.valueOf(cmd.docType()),
            cmd.docNumber(), cmd.issuer(),
            cmd.issueDate(), cmd.expiryDate(),
            kernelFile.secureUrl(),  // URL du Kernel
            null, cmd.notes(), null, null
        );
        return documentPort.saveVehicleDoc(doc);
    });
}
\end{lstlisting}

% ─────────────────────────────────────────────────────────────────────────────
\chapter{Architecture d'Intégration FleetMan — Kernel}
% ─────────────────────────────────────────────────────────────────────────────

\section{Vue d'Ensemble de l'Architecture Cible}

L'architecture cible positionne FleetMan comme un \textbf{backend consommateur} du Kernel RT-Comops. FleetMan conserve sa logique métier propre (gestion de flotte, trajets, géorepérage, opérations terrain) tout en déléguant les fonctions transversales au Kernel.

\begin{figure}[h]
\centering
\begin{tikzpicture}[
  kernel/.style={rectangle, rounded corners=8pt, minimum width=3.5cm,
                 minimum height=1cm, draw=kernelcolor, fill=kernellight,
                 font=\small\bfseries, text=kernelcolor},
  fleet/.style={rectangle, rounded corners=8pt, minimum width=3.5cm,
                minimum height=1cm, draw=fleetblue, fill=fleetblue!10,
                font=\small\bfseries, text=fleetblue},
  arrow/.style={->, >=Stealth, thick},
  node distance=1.2cm and 2.5cm
]
  % Kernel
  \node[kernel] (kcore) at (0,6) {kernel-core};
  \node[kernel] (auth) at (0,4.5) {auth-core};
  \node[kernel] (roles) at (0,3) {roles-core};
  \node[kernel] (org) at (0,1.5) {organization-core};
  \node[kernel] (res) at (0,0) {resource-core};
  \node[kernel] (file) at (0,-1.5) {file-core};
  \node[kernel] (tp) at (0,-3) {tp-core};
  \node[kernel] (settings) at (0,-4.5) {settings-core};

  % FleetMan
  \node[fleet] (vehicles) at (7,4.5) {VehicleService};
  \node[fleet] (drivers) at (7,3) {DriverService};
  \node[fleet] (fleets) at (7,1.5) {FleetService};
  \node[fleet] (docs) at (7,0) {DocumentService};
  \node[fleet] (missions) at (7,-1.5) {MissionService};
  \node[fleet] (trips) at (7,-3) {TripService};
  \node[fleet] (ops) at (7,-4.5) {OperationsService};

  % Arrows
  \draw[arrow, color=kernelcolor] (kcore) -- node[above, font=\tiny] {JWT/Audit} (7,6);
  \draw[arrow, color=kernelcolor] (auth) -- (vehicles);
  \draw[arrow, color=kernelcolor] (roles) -- (drivers);
  \draw[arrow, color=kernelcolor] (org) -- (fleets);
  \draw[arrow, color=kernelcolor] (res) -- (docs);
  \draw[arrow, color=kernelcolor] (file) -- (missions);
  \draw[arrow, color=kernelcolor] (tp) -- (trips);
  \draw[arrow, color=kernelcolor] (settings) -- (ops);

  % Labels
  \node[font=\large\bfseries, color=kernelcolor] at (0,7.2) {Kernel RT-Comops};
  \node[font=\large\bfseries, color=fleetblue] at (7,7.2) {FleetMan Backend};
\end{tikzpicture}
\caption{Architecture d'intégration FleetMan — Kernel RT-Comops}
\end{figure}

\section{Couche d'Intégration : ksm-kernel-client}

Le Kernel fournit un module client commun \texttt{ksm-kernel-client} qui centralise :

\begin{itemize}[leftmargin=2cm]
  \item L'appel HTTP vers le Kernel avec injection automatique des headers
  \item La propagation du contexte (\texttt{X-Tenant-Id}, \texttt{X-Organization-Id})
  \item Le relais du bearer token utilisateur
  \item La gestion des erreurs Kernel (\texttt{CLIENT\_APPLICATION\_SERVICE\_NOT\_ALLOWED}, etc.)
\end{itemize}

\begin{lstlisting}[style=javaStyle, caption={Configuration du ksm-kernel-client dans FleetMan}]
// pom.xml — Ajout de la dependance
<dependency>
    <groupId>com.rt-comops</groupId>
    <artifactId>ksm-kernel-client</artifactId>
    <version>3.0.0</version>
</dependency>

// KernelClientConfig.java
@Configuration
public class KernelClientConfig {

    @Bean
    public KernelClient kernelClient(
            @Value("${application.kernel.url}") String kernelUrl,
            @Value("${application.kernel.client-id}") String clientId,
            @Value("${application.kernel.api-key}") String apiKey) {

        return KernelClient.builder()
                .baseUrl(kernelUrl)
                .clientId(clientId)
                .apiKey(apiKey)
                .build();
    }
}
\end{lstlisting}

\section{Gestion des Erreurs Kernel dans FleetMan}

FleetMan doit traiter les refus fonctionnels du Kernel comme des signaux de contrat :

\begin{longtable}{p{7cm}p{6cm}}
\toprule
\textbf{Code Erreur Kernel} & \textbf{Traitement FleetMan} \\
\midrule
\endhead
\texttt{CLIENT\_APPLICATION\_SERVICE\_NOT\_ALLOWED} & Log critique + alerte admin \\
\texttt{TENANT\_REQUEST\_QUOTA\_EXCEEDED} & HTTP 429 + message utilisateur \\
\texttt{ORGANIZATION\_CONTEXT\_REQUIRED} & HTTP 400 + demande contexte \\
\texttt{ORGANIZATION\_SERVICE\_NOT\_SUBSCRIBED} & HTTP 403 + message abonnement \\
\texttt{ORGANIZATION\_SERVICE\_QUOTA\_EXCEEDED} & HTTP 429 + message quota \\
Refus RBAC & HTTP 403 + message permissions \\
\bottomrule
\caption{Traitement des erreurs Kernel dans FleetMan}
\end{longtable}

\begin{lstlisting}[style=javaStyle, caption={Handler d'erreurs Kernel dans GlobalExceptionHandler}]
@ExceptionHandler(KernelException.class)
public ProblemDetail handleKernelException(KernelException ex) {
    log.error("Erreur Kernel [{}]: {}", ex.getKernelCode(), ex.getMessage());

    return switch (ex.getKernelCode()) {
        case "TENANT_REQUEST_QUOTA_EXCEEDED",
             "ORGANIZATION_SERVICE_QUOTA_EXCEEDED" ->
            ProblemDetail.forStatusAndDetail(HttpStatus.TOO_MANY_REQUESTS,
                "Quota de service atteint. Veuillez reessayer plus tard.");

        case "ORGANIZATION_SERVICE_NOT_SUBSCRIBED" ->
            ProblemDetail.forStatusAndDetail(HttpStatus.FORBIDDEN,
                "Votre organisation n'est pas abonnee a ce service.");

        case "CLIENT_APPLICATION_SERVICE_NOT_ALLOWED" ->
            ProblemDetail.forStatusAndDetail(HttpStatus.FORBIDDEN,
                "Service non autorise pour cette application.");

        default ->
            ProblemDetail.forStatusAndDetail(HttpStatus.BAD_GATEWAY,
                "Erreur du service central: " + ex.getMessage());
    };
}
\end{lstlisting>

% ─────────────────────────────────────────────────────────────────────────────
\chapter{Plan d'Implémentation de l'Intégration}
% ─────────────────────────────────────────────────────────────────────────────

\section{Principes Directeurs}

L'intégration du Kernel dans FleetMan suit les principes suivants :

\begin{enumerate}[leftmargin=2cm]
  \item \textbf{Migration progressive} : FleetMan continue de fonctionner pendant la migration. Les modules locaux sont remplacés un par un.
  \item \textbf{Anti-Corruption Layer} : des adapters dédiés isolent FleetMan des changements d'API du Kernel.
  \item \textbf{Données dénormalisées} : FleetMan conserve des copies locales des données Kernel pour les performances (nom du conducteur, immatriculation du véhicule).
  \item \textbf{Synchronisation par événements} : les changements dans le Kernel sont propagés à FleetMan via Kafka (outbox).
  \item \textbf{Fallback gracieux} : si le Kernel est indisponible, FleetMan dégrade gracieusement (mode dégradé).
\end{enumerate}

\section{Phase 0 — Préparation (Semaine 1-2)}

\begin{tcolorbox}[colback=fleetdark!4, colframe=fleetdark, title={\bfseries\color{white} Phase 0 — Enregistrement et Configuration}]
\textbf{Objectif :} Enregistrer FleetMan comme ClientApplication dans le Kernel et configurer la connexion.

\textbf{Tâches :}
\begin{itemize}[leftmargin=1.5cm]
  \item Créer le tenant FleetMan dans le Kernel (via l'admin Kernel)
  \item Créer la ClientApplication \texttt{fleet-management-backend}
  \item Configurer les \texttt{allowedServiceCodes} : ORGANIZATION, RESOURCE, COMMERCIAL, SETTINGS, FILE
  \item Stocker le secret dans les variables d'environnement
  \item Ajouter la dépendance \texttt{ksm-kernel-client} dans \texttt{pom.xml}
  \item Créer \texttt{KernelClientConfig} et \texttt{KernelWebClientAdapter}
  \item Tester la connexion via l'endpoint \texttt{/kernel/health}
\end{itemize}

\textbf{Livrable :} FleetMan connecté au Kernel, headers propagés correctement.
\end{tcolorbox}

\section{Phase 1 — Intégration auth-core et roles-core (Semaine 3-5)}

\begin{tcolorbox}[colback=fleetblue!5, colframe=fleetblue, title={\bfseries\color{white} Phase 1 — Authentification et RBAC Centralisés}]
\textbf{Objectif :} Remplacer l'authentification locale par l'auth-core du Kernel.

\textbf{Tâches :}
\begin{itemize}[leftmargin=1.5cm]
  \item Configurer la validation JWT RS256 via JWKS du Kernel
  \item Adapter \texttt{SecurityConfig} pour utiliser \texttt{oauth2ResourceServer}
  \item Migrer \texttt{JwtAuthenticationManager} vers le validateur Kernel
  \item Provisionner les templates de rôles FleetMan dans le Kernel
  \item Adapter les \texttt{@PreAuthorize} pour utiliser les permissions granulaires
  \item Migrer les utilisateurs existants vers le Kernel (script de migration)
  \item Tester les flows login, refresh, logout
\end{itemize}

\textbf{Livrable :} Authentification 100\% Kernel, rôles RBAC granulaires.
\end{tcolorbox}

\section{Phase 2 — Intégration organization-core (Semaine 6-7)}

\begin{tcolorbox}[colback=kernellight, colframe=kernelcolor, title={\bfseries\color{white} Phase 2 — Flottes comme Organisations}]
\textbf{Objectif :} Synchroniser les flottes FleetMan avec les organisations Kernel.

\textbf{Tâches :}
\begin{itemize}[leftmargin=1.5cm]
  \item Créer \texttt{KernelOrganizationAdapter} implémentant un port sortant
  \item Modifier \texttt{FleetService.createFleet()} pour créer une organisation Kernel
  \item Ajouter le champ \texttt{kernel\_organization\_id} à \texttt{fleet.fleets}
  \item Créer \texttt{OrganizationEventConsumer} pour synchroniser les suspensions
  \item Propager \texttt{X-Organization-Id} dans tous les appels FleetMan
  \item Migrer les flottes existantes vers des organisations Kernel
\end{itemize}

\textbf{Livrable :} Flottes synchronisées avec le Kernel, multi-tenancy opérationnel.
\end{tcolorbox}

\section{Phase 3 — Intégration resource-core (Semaine 8-10)}

\begin{tcolorbox}[colback=fleetblue!5, colframe=fleetblue, title={\bfseries\color{white} Phase 3 — Véhicules comme Ressources}]
\textbf{Objectif :} Remplacer le service véhicule externe par resource-core.

\textbf{Tâches :}
\begin{itemize}[leftmargin=1.5cm]
  \item Créer \texttt{KernelResourceAdapter} remplaçant \texttt{VehicleApiAdapter}
  \item Modifier \texttt{VehicleService} pour déléguer la création au Kernel
  \item Synchroniser les affectations conducteur-véhicule avec resource-core
  \item Intégrer la mise à jour GPS via le port resource-core
  \item Synchroniser les passages en maintenance avec resource-core
  \item Migrer les véhicules existants vers des ressources Kernel
\end{itemize}

\textbf{Livrable :} Véhicules gérés par resource-core, GPS synchronisé.
\end{tcolorbox}

\section{Phase 4 — Intégration actor-core et tp-core (Semaine 11-12)}

\begin{tcolorbox}[colback=kernellight, colframe=kernelcolor, title={\bfseries\color{white} Phase 4 — Acteurs et Tiers}]
\textbf{Objectif :} Synchroniser conducteurs et clients avec le Kernel.

\textbf{Tâches :}
\begin{itemize}[leftmargin=1.5cm]
  \item Créer \texttt{KernelActorAdapter} pour la synchronisation des conducteurs
  \item Modifier \texttt{DriverService} pour créer un acteur Kernel à l'enregistrement
  \item Adapter le Module 7 (Clients) pour utiliser tp-core
  \item Créer \texttt{KernelThirdPartyAdapter}
  \item Migrer les conducteurs et clients existants
\end{itemize}

\textbf{Livrable :} Conducteurs et clients synchronisés avec le Kernel.
\end{tcolorbox}

\section{Phase 5 — Intégration file-core et settings-core (Semaine 13-14)}

\begin{tcolorbox}[colback=fleetblue!5, colframe=fleetblue, title={\bfseries\color{white} Phase 5 — Fichiers et Paramètres}]
\textbf{Objectif :} Centraliser les fichiers et paramètres dans le Kernel.

\textbf{Tâches :}
\begin{itemize}[leftmargin=1.5cm]
  \item Créer \texttt{KernelFileAdapter} pour le téléversement des documents légaux
  \item Modifier \texttt{DocumentService} pour utiliser file-core
  \item Créer \texttt{KernelSettingsAdapter} pour la numérotation des missions
  \item Modifier \texttt{MissionService} pour utiliser settings-core
  \item Migrer les fichiers existants vers le Kernel
\end{itemize}

\textbf{Livrable :} Documents légaux et numérotation gérés par le Kernel.
\end{tcolorbox}

\section{Phase 6 — Intégration common-core et Finalisation (Semaine 15-16)}

\begin{tcolorbox}[colback=kernellight, colframe=kernelcolor, title={\bfseries\color{white} Phase 6 — Types Partagés et Finalisation}]
\textbf{Objectif :} Aligner les types FleetMan sur le common-core et finaliser l'intégration.

\textbf{Tâches :}
\begin{itemize}[leftmargin=1.5cm]
  \item Aligner \texttt{PageResponse<T>} sur \texttt{PageResult<T>} du Kernel
  \item Adopter \texttt{ApiResponse<T>} pour toutes les réponses REST
  \item Ajouter \texttt{tenantId} sur toutes les entités FleetMan
  \item Intégrer les adresses polymorphiques pour les flottes et dépôts
  \item Tests d'intégration end-to-end
  \item Documentation de l'API mise à jour
\end{itemize}

\textbf{Livrable :} Intégration complète, tests validés, documentation à jour.
\end{tcolorbox}

\section{Tableau Récapitulatif du Plan}

\begin{longtable}{p{1cm}p{4cm}p{3cm}p{2cm}p{3cm}}
\toprule
\textbf{Phase} & \textbf{Modules Kernel} & \textbf{Semaines} & \textbf{Effort} & \textbf{Dépendances} \\
\midrule
\endhead
0 & Enregistrement & 1-2 & S & Aucune \\
1 & auth-core, roles-core & 3-5 & M & Phase 0 \\
2 & organization-core & 6-7 & M & Phase 1 \\
3 & resource-core & 8-10 & L & Phase 2 \\
4 & actor-core, tp-core & 11-12 & M & Phase 2 \\
5 & file-core, settings-core & 13-14 & M & Phase 3, 4 \\
6 & common-core, finalisation & 15-16 & S & Phase 1-5 \\
\bottomrule
\caption{Plan d'implémentation de l'intégration Kernel (16 semaines)}
\end{longtable}
% ─────────────────────────────────────────────────────────────────────────────
\chapter{Impact sur l'Architecture Existante de FleetMan}
% ─────────────────────────────────────────────────────────────────────────────

\section{Composants à Remplacer}

L'intégration du Kernel implique le remplacement de plusieurs composants
existants de FleetMan par des adapters vers le Kernel :

\begin{longtable}{p{5.5cm}p{5.5cm}p{2cm}}
\toprule
\textbf{Composant FleetMan Actuel} & \textbf{Remplacement Kernel} & \textbf{Phase} \\
\midrule
\endhead
\texttt{RemoteAuthAdapter} & \texttt{KernelAuthAdapter} & 1 \\
\texttt{FakeAuthAdapter} & Supprimé & 1 \\
\texttt{JwtAuthenticationManager} & Validateur JWKS Kernel & 1 \\
\texttt{BearerTokenServerAuthenticationConverter} & Convertisseur Kernel & 1 \\
\texttt{SystemUserInitializer} & Provisionnement Kernel & 1 \\
\texttt{VehicleApiAdapter} & \texttt{KernelResourceAdapter} & 3 \\
\texttt{VehicleApiClient} & \texttt{KernelWebClient} & 3 \\
\texttt{AuthApiClient} & \texttt{KernelAuthClient} & 1 \\
\bottomrule
\caption{Composants FleetMan remplacés par des adapters Kernel}
\end{longtable}

\section{Composants à Enrichir}

Ces composants existants sont conservés mais enrichis avec des appels Kernel :

\begin{longtable}{p{5cm}p{8cm}}
\toprule
\textbf{Composant} & \textbf{Enrichissement} \\
\midrule
\endhead
\texttt{FleetService} & Création d'organisation Kernel à la création de flotte \\
\texttt{DriverService} & Création d'acteur Kernel à l'enregistrement du conducteur \\
\texttt{VehicleService} & Enregistrement ressource Kernel + synchronisation GPS \\
\texttt{MaintenanceService} & Synchronisation état ressource Kernel \\
\texttt{DocumentService} & Téléversement fichiers via file-core Kernel \\
\texttt{MissionService} & Numérotation via settings-core + tiers via tp-core \\
\texttt{GlobalExceptionHandler} & Gestion des erreurs Kernel \\
\texttt{SecurityConfig} & Validation JWT RS256 via JWKS Kernel \\
\bottomrule
\caption{Composants FleetMan enrichis avec des appels Kernel}
\end{longtable}

\section{Nouveaux Composants à Créer}

\begin{longtable}{p{6cm}p{7cm}}
\toprule
\textbf{Nouveau Composant} & \textbf{Rôle} \\
\midrule
\endhead
\texttt{KernelClientConfig} & Configuration du client Kernel \\
\texttt{KernelWebClientAdapter} & Injection des headers Kernel \\
\texttt{KernelAuthAdapter} & Authentification via auth-core \\
\texttt{KernelResourceAdapter} & Gestion ressources via resource-core \\
\texttt{KernelOrganizationAdapter} & Synchronisation organisations \\
\texttt{KernelActorAdapter} & Synchronisation acteurs \\
\texttt{KernelThirdPartyAdapter} & Gestion tiers via tp-core \\
\texttt{KernelFileAdapter} & Téléversement fichiers via file-core \\
\texttt{KernelSettingsAdapter} & Paramètres et numérotation \\
\texttt{OrganizationEventConsumer} & Consommateur événements Kafka Kernel \\
\texttt{ResourceEventConsumer} & Consommateur événements ressources \\
\texttt{KernelException} & Exception typée pour les erreurs Kernel \\
\bottomrule
\caption{Nouveaux composants à créer pour l'intégration Kernel}
\end{longtable}

\section{Impact sur la Base de Données}

L'intégration nécessite l'ajout de colonnes de référence vers le Kernel :

\begin{lstlisting}[style=yamlStyle, caption={Migration SQL pour les références Kernel}]
-- Migration 020-add-kernel-references.sql

-- Flottes : reference vers l'organisation Kernel
ALTER TABLE fleet.fleets
    ADD COLUMN IF NOT EXISTS kernel_organization_id UUID,
    ADD COLUMN IF NOT EXISTS kernel_tenant_id UUID;

-- Vehicules : reference vers la ressource Kernel
ALTER TABLE fleet.vehicles
    ADD COLUMN IF NOT EXISTS kernel_resource_id UUID;

-- Conducteurs : reference vers l'acteur Kernel
ALTER TABLE fleet.drivers
    ADD COLUMN IF NOT EXISTS kernel_actor_id UUID;

-- Missions : reference vers le tiers Kernel
ALTER TABLE fleet.missions
    ADD COLUMN IF NOT EXISTS kernel_third_party_id UUID;

-- Documents : reference vers le fichier Kernel
ALTER TABLE fleet.vehicle_documents
    ADD COLUMN IF NOT EXISTS kernel_file_id UUID;
ALTER TABLE fleet.driver_documents
    ADD COLUMN IF NOT EXISTS kernel_file_id UUID;

-- Index pour les jointures
CREATE INDEX IF NOT EXISTS idx_fleets_kernel_org
    ON fleet.fleets(kernel_organization_id);
CREATE INDEX IF NOT EXISTS idx_vehicles_kernel_res
    ON fleet.vehicles(kernel_resource_id);
CREATE INDEX IF NOT EXISTS idx_drivers_kernel_actor
    ON fleet.drivers(kernel_actor_id);
\end{lstlisting}

% ─────────────────────────────────────────────────────────────────────────────
\chapter{Patterns d'Intégration Détaillés}
% ─────────────────────────────────────────────────────────────────────────────

\section{Pattern 1 — Anti-Corruption Layer (ACL)}

Chaque module Kernel est encapsulé dans un adapter dédié qui isole FleetMan
des changements d'API du Kernel. Ce pattern garantit que si l'API Kernel évolue,
seul l'adapter doit être modifié.

\begin{lstlisting}[style=javaStyle, caption={Pattern ACL — KernelResourceAdapter}]
/**
 * Anti-Corruption Layer vers resource-core du Kernel.
 * Isole FleetMan des changements d'API du Kernel.
 * Implémente ExternalVehiclePort (port sortant existant).
 */
@Component
@RequiredArgsConstructor
@Slf4j
public class KernelResourceAdapter implements ExternalVehiclePort {

    private final KernelWebClientAdapter kernelClient;
    private final WebClient webClient;

    @Override
    public Mono<String> registerVehicle(VehicleRequest request,
                                         String organizationId,
                                         String token) {
        // Traduction du modele FleetMan vers le modele Kernel
        Map<String, Object> kernelRequest = Map.of(
            "resourceType", "VEHICLE",
            "name", request.brand() + " " + request.model(),
            "serialNumber", request.vehicleSerialNumber(),
            "attributes", Map.of(
                "licensePlate", request.licensePlate(),
                "color", request.color(),
                "year", request.manufacturingYear()
            )
        );

        return kernelClient.withKernelHeaders(
            webClient.post().uri("/api/v1/resources")
                     .bodyValue(kernelRequest),
            organizationId, token
        )
        .retrieve()
        .onStatus(status -> status.is4xxClientError(),
            resp -> resp.bodyToMono(String.class)
                .flatMap(body -> Mono.error(
                    new KernelException("RESOURCE_CREATION_FAILED", body))))
        .bodyToMono(KernelResourceResponse.class)
        .map(KernelResourceResponse::id)
        .doOnSuccess(id -> log.info(
            "Vehicule enregistre dans Kernel resource-core: {}", id));
    }

    // DTO interne de reponse Kernel (isole du domaine FleetMan)
    private record KernelResourceResponse(String id, String status) {}
}
\end{lstlisting}

\section{Pattern 2 — Synchronisation par Événements Kafka}

Les changements dans le Kernel sont propagés à FleetMan via Kafka.
FleetMan consomme les événements du topic \texttt{iwm.events.business}
et met à jour ses projections locales.

\begin{lstlisting}[style=javaStyle, caption={Consommateur d'événements Kernel}]
@Component
@RequiredArgsConstructor
@Slf4j
public class KernelEventConsumer {

    private final FleetRepositoryPort fleetPort;
    private final VehiclePersistencePort vehiclePort;
    private final DriverPersistencePort driverPort;

    @KafkaListener(
        topics = "iwm.events.business",
        groupId = "fleet-management-group",
        containerFactory = "kernelEventListenerFactory"
    )
    public void onKernelEvent(KernelBusinessEvent event) {
        log.debug("Evenement Kernel recu: type={}, entityId={}",
                  event.type(), event.entityId());

        switch (event.type()) {
            case "ORGANIZATION_SUSPENDED" ->
                handleOrganizationSuspended(event.entityId());
            case "RESOURCE_STATUS_CHANGED" ->
                handleResourceStatusChanged(event.entityId(),
                                            event.payload());
            case "ACTOR_DEACTIVATED" ->
                handleActorDeactivated(event.entityId());
            default ->
                log.debug("Evenement Kernel ignore: {}", event.type());
        }
    }

    private void handleOrganizationSuspended(String orgId) {
        fleetPort.findByKernelOrganizationId(UUID.fromString(orgId))
            .flatMap(fleet -> fleetPort.suspend(fleet.id()))
            .doOnSuccess(f -> log.info(
                "Flotte suspendue suite a suspension organisation Kernel: {}",
                orgId))
            .subscribe();
    }

    private void handleResourceStatusChanged(String resourceId,
                                              Map<String, Object> payload) {
        String newStatus = (String) payload.get("status");
        vehiclePort.findByKernelResourceId(UUID.fromString(resourceId))
            .flatMap(v -> vehiclePort.updateStatus(v.id(), newStatus))
            .subscribe();
    }

    private void handleActorDeactivated(String actorId) {
        driverPort.findByKernelActorId(UUID.fromString(actorId))
            .flatMap(d -> driverPort.deactivate(d.userId()))
            .subscribe();
    }
}
\end{lstlisting}

\section{Pattern 3 — Données Dénormalisées pour la Performance}

FleetMan conserve des copies locales des données Kernel pour éviter
des appels HTTP synchrones sur chaque requête. Ces copies sont
maintenues à jour via les événements Kafka.

\begin{lstlisting}[style=javaStyle, caption={Strategie de denormalisation}]
// Dans VehicleService.createIndependentVehicle()
// Apres enregistrement dans le Kernel :

Vehicle vehicle = new Vehicle(
    null,
    fleetId,
    managerId,
    null,
    vehicleTypeId,
    request.licensePlate(),
    request.vehicleSerialNumber(),
    request.brand(),
    request.model(),
    // ... autres champs locaux ...
    kernelResourceId  // Reference vers le Kernel
);

// La donnee est stockee localement ET referencee dans le Kernel.
// Les lectures utilisent la copie locale (performance).
// Les ecritures passent par le Kernel (coherence).
\end{lstlisting}

\section{Pattern 4 — Fallback Gracieux}

Si le Kernel est temporairement indisponible, FleetMan doit dégrader
gracieusement sans bloquer les opérations critiques.

\begin{lstlisting}[style=javaStyle, caption={Circuit Breaker vers le Kernel}]
@Component
@RequiredArgsConstructor
public class KernelResourceAdapter implements ExternalVehiclePort {

    private final CircuitBreakerFactory circuitBreakerFactory;

    @Override
    public Mono<String> registerVehicle(VehicleRequest request,
                                         String orgId, String token) {
        CircuitBreaker cb = circuitBreakerFactory
            .create("kernel-resource");

        return cb.run(
            // Appel normal vers le Kernel
            doRegisterVehicle(request, orgId, token),
            // Fallback : enregistrement local sans Kernel
            throwable -> {
                log.warn("Kernel indisponible, mode degrade: {}",
                         throwable.getMessage());
                // Retourner un ID temporaire local
                return Mono.just("LOCAL-" + UUID.randomUUID());
            }
        );
    }
}
\end{lstlisting>

% ─────────────────────────────────────────────────────────────────────────────
\chapter{Impacts sur la Sécurité et la Conformité}
% ─────────────────────────────────────────────────────────────────────────────

\section{Modèle de Sécurité Cible}

Avec l'intégration du Kernel, le modèle de sécurité de FleetMan évolue
vers un modèle à trois couches :

\begin{enumerate}[leftmargin=2cm]
  \item \textbf{Couche Kernel} : authentification, validation JWT RS256,
        vérification ClientApplication, quotas
  \item \textbf{Couche FleetMan} : autorisation RBAC via permissions Kernel,
        cloisonnement par flotte (organizationId)
  \item \textbf{Couche Données} : filtrage systématique par \texttt{tenantId}
        et \texttt{managerId} dans toutes les requêtes
\end{enumerate}

\section{Cloisonnement Multi-Tenant}

Chaque entreprise de transport devient un tenant distinct dans le Kernel.
Le cloisonnement est garanti par :

\begin{itemize}[leftmargin=2cm]
  \item Le \texttt{tenantId} présent dans le JWT et propagé dans toute la
        chaîne réactive
  \item Le filtre \texttt{TenantWebFilter} du Kernel qui injecte le contexte
  \item Les requêtes SQL de FleetMan qui filtrent systématiquement par
        \texttt{manager\_id} (lié à l'\texttt{organizationId} Kernel)
  \item L'impossibilité pour un gestionnaire d'une flotte d'accéder aux
        données d'une autre flotte (même tenant, organisations différentes)
\end{itemize}

\section{Audit et Traçabilité}

L'intégration du Kernel enrichit considérablement la traçabilité de FleetMan :

\begin{longtable}{p{5cm}p{8cm}}
\toprule
\textbf{Opération FleetMan} & \textbf{Trace dans le Kernel} \\
\midrule
\endhead
Création d'un véhicule & \texttt{AuditTrail} Kernel + événement outbox \\
Enregistrement d'un conducteur & \texttt{AuditTrail} Kernel (acteur créé) \\
Déclaration d'un incident & Événement Kafka \texttt{INCIDENT\_REPORTED} \\
Téléversement d'un document & \texttt{AuditTrail} Kernel (file-core) \\
Connexion d'un utilisateur & \texttt{AuditTrail} Kernel (auth-core) \\
Modification d'un rôle & \texttt{AuditTrail} Kernel (roles-core) \\
\bottomrule
\caption{Traçabilité enrichie par le Kernel}
\end{longtable}

\section{Gestion des Secrets}

Les secrets de la ClientApplication FleetMan doivent être gérés
selon les bonnes pratiques :

\begin{lstlisting}[style=yamlStyle, caption={Gestion des secrets en production}]
# Variables d'environnement (jamais en dur dans le code)
FLEET_KERNEL_API_KEY=<secret-genere-par-le-kernel>
FLEET_KERNEL_CLIENT_ID=fleet-management-backend
FLEET_TENANT_ID=<uuid-du-tenant-fleetman>

# En production : utiliser un gestionnaire de secrets
# (HashiCorp Vault, AWS Secrets Manager, etc.)
# Le secret doit etre pivote regulierement via l'API Kernel
\end{lstlisting}

% ─────────────────────────────────────────────────────────────────────────────
\chapter{Bénéfices et Risques de l'Intégration}
% ─────────────────────────────────────────────────────────────────────────────

\section{Bénéfices Attendus}

\subsection{Bénéfices Fonctionnels}

\begin{longtable}{p{4cm}p{9cm}}
\toprule
\textbf{Domaine} & \textbf{Bénéfice} \\
\midrule
\endhead
Authentification & Login multi-contexte, OTP, MFA, réinitialisation mot de passe \\
Gestion des droits & Permissions granulaires, scopes organisationnels \\
Multi-tenancy & Isolation complète entre entreprises de transport \\
Traçabilité & Audit immuable sur toutes les opérations critiques \\
Fichiers & Stockage centralisé, gouvernance documentaire, expiration \\
Clients & Cycle de vie complet prospect → client, recherche plein-texte \\
Numérotation & Séquences documentaires standardisées pour missions/factures \\
Véhicules & Cycle de vie unifié, localisation GPS centralisée \\
\bottomrule
\caption{Bénéfices fonctionnels de l'intégration Kernel}
\end{longtable}

\subsection{Bénéfices Techniques}

\begin{itemize}[leftmargin=2cm]
  \item \textbf{Réduction du code} : suppression de ~2000 lignes de code
        d'authentification et de gestion des droits locaux
  \item \textbf{Cohérence} : même modèle de données que les autres modules
        du projet RT-Comops
  \item \textbf{Évolutivité} : les améliorations du Kernel bénéficient
        automatiquement à FleetMan
  \item \textbf{Observabilité} : métriques Prometheus unifiées sur tout
        l'écosystème
  \item \textbf{Sécurité} : JWT RS256 standardisé, rotation des secrets
        centralisée
\end{itemize}

\section{Risques et Mitigations}

\begin{longtable}{p{4cm}p{5cm}p{4cm}}
\toprule
\textbf{Risque} & \textbf{Impact} & \textbf{Mitigation} \\
\midrule
\endhead
Indisponibilité Kernel & Blocage des opérations & Circuit Breaker + mode dégradé \\
Latence réseau & Ralentissement des APIs & Cache Redis local + dénormalisation \\
Rupture d'API Kernel & Régression fonctionnelle & Anti-Corruption Layer (ACL) \\
Migration des données & Perte de données & Migration progressive + rollback \\
Quota dépassé & Refus de service & Monitoring + alertes préventives \\
Secret compromis & Accès non autorisé & Rotation régulière + audit \\
\bottomrule
\caption{Risques et mitigations de l'intégration Kernel}
\end{longtable}

\section{Métriques de Succès}

L'intégration sera considérée réussie si :

\begin{itemize}[leftmargin=2cm]
  \item \textbf{Fonctionnel} : 100\% des flows d'authentification passent
        par le Kernel
  \item \textbf{Performance} : latence P95 $\leq$ 200ms sur les endpoints
        FleetMan (incluant les appels Kernel)
  \item \textbf{Disponibilité} : FleetMan reste opérationnel en mode dégradé
        si le Kernel est indisponible $\leq$ 30 secondes
  \item \textbf{Sécurité} : 0 accès cross-tenant détecté en audit
  \item \textbf{Cohérence} : 100\% des véhicules FleetMan ont un
        \texttt{kernel\_resource\_id} valide
\end{itemize}

% ─────────────────────────────────────────────────────────────────────────────
\chapter{Spécifications Techniques des Adapters}
% ─────────────────────────────────────────────────────────────────────────────

\section{Structure du Package d'Intégration}

Tous les adapters Kernel sont regroupés dans un package dédié :

\begin{lstlisting}[style=yamlStyle, caption={Structure du package d'intégration Kernel}]
src/main/java/com/yowyob/fleet/
  infrastructure/
    adapters/
      outbound/
        kernel/                          # Nouveau package
          auth/
            KernelAuthAdapter.java       # auth-core
          resource/
            KernelResourceAdapter.java   # resource-core
          organization/
            KernelOrganizationAdapter.java # organization-core
          actor/
            KernelActorAdapter.java      # actor-core
          thirdparty/
            KernelThirdPartyAdapter.java # tp-core
          file/
            KernelFileAdapter.java       # file-core
          settings/
            KernelSettingsAdapter.java   # settings-core
          client/
            KernelWebClientAdapter.java  # Client commun
            KernelClientConfig.java      # Configuration
          dto/
            KernelResourceResponse.java  # DTOs internes
            KernelOrganizationResponse.java
            KernelActorResponse.java
            KernelFileResponse.java
          exception/
            KernelException.java         # Exception typee
      inbound/
        messaging/
          KernelEventConsumer.java       # Consommateur Kafka
\end{lstlisting}

\section{Interface Commune des Adapters Kernel}

Tous les adapters Kernel implémentent une interface commune
pour la gestion du contexte :

\begin{lstlisting}[style=javaStyle, caption={Interface commune KernelAdapter}]
/**
 * Interface commune pour tous les adapters Kernel.
 * Garantit la propagation coherente du contexte.
 */
public interface KernelAdapter {

    /**
     * Construit les headers Kernel pour un appel authentifie.
     * @param organizationId ID de l'organisation (X-Organization-Id)
     * @param bearerToken    Token JWT de l'utilisateur
     */
    default WebClient.RequestHeadersSpec<?> withContext(
            WebClient.RequestHeadersSpec<?> spec,
            String organizationId,
            String bearerToken) {
        return spec
            .header("X-Client-Id", getClientId())
            .header("X-Api-Key", getApiKey())
            .header("X-Tenant-Id", getTenantId())
            .header("X-Organization-Id", organizationId)
            .header("Authorization", "Bearer " + bearerToken);
    }

    String getClientId();
    String getApiKey();
    String getTenantId();
}
\end{lstlisting}

\section{Configuration application.yml Enrichie}

\begin{lstlisting}[style=yamlStyle, caption={Configuration complète de l'intégration Kernel}]
application:
  # Configuration Kernel RT-Comops
  kernel:
    url: ${KERNEL_URL:https://kernel.rt-comops.com}
    client-id: ${FLEET_KERNEL_CLIENT_ID:fleet-management-backend}
    api-key: ${FLEET_KERNEL_API_KEY}
    tenant-id: ${FLEET_TENANT_ID}
    timeout:
      connect: 5s
      read: 10s
    retry:
      max-attempts: 3
      backoff: 1s
    circuit-breaker:
      failure-rate-threshold: 50
      wait-duration: 30s
    allowed-services:
      - ORGANIZATION
      - RESOURCE
      - COMMERCIAL
      - SETTINGS
      - FILE

  # Topics Kafka Kernel
  kafka:
    topics:
      kernel-events: iwm.events.business
      operation-events: operation-events-topic
      geofence-alerts: geofence-alerts-topic
\end{lstlisting}

% ─────────────────────────────────────────────────────────────────────────────
\chapter{Synthèse et Recommandations}
% ─────────────────────────────────────────────────────────────────────────────

\section{Synthèse des Modules Kernel Retenus}

\begin{longtable}{p{3.5cm}p{3cm}p{7cm}}
\toprule
\textbf{Module Kernel} & \textbf{Priorité} & \textbf{Valeur pour FleetMan} \\
\midrule
\endhead
\texttt{kernel-core} & \textcolor{fleeterror}{\textbf{Critique}} &
  Fondation obligatoire : multi-tenancy, JWT, audit, outbox \\
\texttt{auth-core} & \textcolor{fleeterror}{\textbf{Critique}} &
  Remplace l'auth locale, apporte MFA et multi-contexte \\
\texttt{roles-core} & \textcolor{fleeterror}{\textbf{Critique}} &
  RBAC granulaire avec cache Redis \\
\texttt{organization-core} & \textcolor{fleetwarning}{\textbf{Haute}} &
  Flotte = Organisation, dépôt = Agence \\
\texttt{resource-core} & \textcolor{fleetwarning}{\textbf{Haute}} &
  Véhicules comme ressources, GPS centralisé \\
\texttt{common-core} & \textcolor{fleetwarning}{\textbf{Haute}} &
  Types partagés, cohérence structurelle \\
\texttt{file-core} & \textcolor{fleetsuccess}{\textbf{Normale}} &
  Documents légaux centralisés \\
\texttt{tp-core} & \textcolor{fleetsuccess}{\textbf{Normale}} &
  Clients des missions comme tiers \\
\texttt{settings-core} & \textcolor{fleetsuccess}{\textbf{Normale}} &
  Numérotation et paramètres métier \\
\texttt{actor-core} & \textcolor{fleetsuccess}{\textbf{Normale}} &
  Conducteurs comme acteurs \\
\bottomrule
\caption{Synthèse des modules Kernel retenus pour FleetMan}
\end{longtable}

\section{Recommandations Finales}

\begin{enumerate}[leftmargin=2cm]
  \item \textbf{Commencer par kernel-core + auth-core} : ces deux modules
        sont les plus structurants et doivent être intégrés en premier.
        Toute la sécurité de FleetMan en dépend.

  \item \textbf{Adopter le pattern ACL systématiquement} : chaque module
        Kernel doit avoir son propre adapter isolé. Ne jamais appeler
        directement les APIs Kernel depuis les services métier.

  \item \textbf{Maintenir les données dénormalisées} : pour les performances,
        FleetMan doit conserver des copies locales des données Kernel
        (nom du conducteur, immatriculation, etc.) et les synchroniser
        via Kafka.

  \item \textbf{Implémenter le Circuit Breaker dès le début} : le Kernel
        est un service externe. FleetMan doit être résilient à son
        indisponibilité.

  \item \textbf{Tester la migration des données} : avant chaque phase,
        tester la migration des données existantes vers le Kernel sur
        un environnement de staging.

  \item \textbf{Documenter les contrats d'intégration} : maintenir un
        document vivant des endpoints Kernel utilisés par FleetMan,
        avec les versions et les changements.

  \item \textbf{Monitorer les quotas} : configurer des alertes Prometheus
        lorsque les quotas Kernel atteignent 80\% pour éviter les
        interruptions de service.
\end{enumerate}

\section{Roadmap Recommandée}

\begin{figure}[h]
\centering
\begin{tikzpicture}[
  phase/.style={rectangle, rounded corners=6pt, minimum width=3cm,
                minimum height=0.8cm, draw=fleetblue, fill=fleetblue!10,
                font=\small\bfseries},
  critical/.style={rectangle, rounded corners=6pt, minimum width=3cm,
                   minimum height=0.8cm, draw=fleeterror, fill=fleeterror!10,
                   font=\small\bfseries},
  arrow/.style={->, >=Stealth, thick, color=fleetdark},
  node distance=0.5cm and 1.5cm
]
  \node[critical] (p0) at (0,0) {Phase 0 : Setup};
  \node[critical] (p1) at (3.5,0) {Phase 1 : Auth};
  \node[phase] (p2) at (7,0) {Phase 2 : Org};
  \node[phase] (p3) at (10.5,0) {Phase 3 : Resource};
  \node[phase] (p4) at (3.5,-2) {Phase 4 : Acteurs};
  \node[phase] (p5) at (7,-2) {Phase 5 : Files};
  \node[phase] (p6) at (10.5,-2) {Phase 6 : Final};

  \draw[arrow] (p0) -- (p1);
  \draw[arrow] (p1) -- (p2);
  \draw[arrow] (p2) -- (p3);
  \draw[arrow] (p1) -- (p4);
  \draw[arrow] (p3) -- (p5);
  \draw[arrow] (p4) -- (p5);
  \draw[arrow] (p5) -- (p6);

  \node[font=\tiny, color=fleetgray] at (0,-0.7) {S1-2};
  \node[font=\tiny, color=fleetgray] at (3.5,-0.7) {S3-5};
  \node[font=\tiny, color=fleetgray] at (7,-0.7) {S6-7};
  \node[font=\tiny, color=fleetgray] at (10.5,-0.7) {S8-10};
  \node[font=\tiny, color=fleetgray] at (3.5,-2.7) {S11-12};
  \node[font=\tiny, color=fleetgray] at (7,-2.7) {S13-14};
  \node[font=\tiny, color=fleetgray] at (10.5,-2.7) {S15-16};
\end{tikzpicture}
\caption{Roadmap d'intégration Kernel dans FleetMan (16 semaines)}
\end{figure}

% ─────────────────────────────────────────────────────────────────────────────
\chapter*{Conclusion}
\addcontentsline{toc}{chapter}{Conclusion}
\thispagestyle{fancy}
% ─────────────────────────────────────────────────────────────────────────────

Ce document a présenté une analyse complète des 20 cores du Kernel RT-Comops
et identifié les 9 modules directement applicables au projet FleetMan.

L'intégration du Kernel dans FleetMan représente une évolution architecturale
majeure qui transforme FleetMan d'une application autonome en un module
métier cohérent de l'écosystème RT-Comops. Les bénéfices sont considérables :

\begin{itemize}[leftmargin=2cm]
  \item \textbf{Sécurité renforcée} : JWT RS256 standardisé, MFA, audit immuable
  \item \textbf{Multi-tenancy réel} : isolation complète entre entreprises de transport
  \item \textbf{Cohérence globale} : même modèle de données que les autres modules
  \item \textbf{Réduction de la dette technique} : suppression du code d'auth local
  \item \textbf{Évolutivité} : bénéfice automatique des améliorations du Kernel
\end{itemize}

Le plan d'implémentation en 6 phases sur 16 semaines garantit une migration
progressive et sécurisée, sans interruption de service. L'adoption du pattern
Anti-Corruption Layer protège FleetMan des évolutions futures du Kernel.

\begin{tcolorbox}[colback=kernellight, colframe=kernelcolor,
                  title={\bfseries\color{white} Message Clé}]
FleetMan n'est pas un simple client HTTP du Kernel. Il est une
\textbf{application consommatrice gouvernée} par le Kernel. La connexion
est valide seulement si l'identité machine, le tenant, l'organisation,
le service, le quota et le JWT utilisateur sont cohérents. Cette gouvernance
est la garantie de la symbiose et de la synchronicité globales du projet
RT-Comops.
\end{tcolorbox}

% ─────────────────────────────────────────────────────────────────────────────
% ANNEXES
% ─────────────────────────────────────────────────────────────────────────────
\appendix

\chapter{Glossaire des Termes Kernel}

\begin{longtable}{p{4cm}p{9cm}}
\toprule
\textbf{Terme} & \textbf{Définition} \\
\midrule
\endhead
\textbf{Tenant} & Unité d'isolation logique de la plateforme (une entreprise) \\
\textbf{ClientApplication} & Identité machine d'un backend consommateur \\
\textbf{allowedServiceCodes} & Liste des services autorisés pour une ClientApplication \\
\textbf{JWT RS256} & Token d'accès signé avec clé RSA 256 bits \\
\textbf{JWKS} & JSON Web Key Set — document exposant les clés publiques \\
\textbf{Outbox} & Table transactionnelle pour la publication fiable d'événements \\
\textbf{AuditTrail} & Piste d'audit immuable (insertion seule) \\
\textbf{PartyRef} & Référence légère inter-cores (partyType + partyId) \\
\textbf{BaseEntity} & Classe abstraite avec UUID, tenantId, timestamps \\
\textbf{PlatformServiceCode} & Énumération des services disponibles dans le Kernel \\
\textbf{ACL} & Anti-Corruption Layer — adapter isolant un domaine externe \\
\textbf{ksm-kernel-client} & Module client commun pour les backends consommateurs \\
\textbf{X-Client-Id} & Header identifiant la ClientApplication \\
\textbf{X-Api-Key} & Header portant le secret de la ClientApplication \\
\textbf{X-Tenant-Id} & Header identifiant le tenant cible \\
\textbf{X-Organization-Id} & Header identifiant l'organisation cible \\
\bottomrule
\caption{Glossaire des termes Kernel}
\end{longtable}

\chapter{Checklist d'Intégration}

\begin{longtable}{p{0.5cm}p{12cm}}
\toprule
& \textbf{Tâche} \\
\midrule
\endhead
$\square$ & Créer le tenant FleetMan dans le Kernel \\
$\square$ & Créer la ClientApplication \texttt{fleet-management-backend} \\
$\square$ & Configurer les \texttt{allowedServiceCodes} \\
$\square$ & Stocker le secret dans les variables d'environnement \\
$\square$ & Ajouter la dépendance \texttt{ksm-kernel-client} \\
$\square$ & Créer \texttt{KernelClientConfig} \\
$\square$ & Configurer la validation JWT RS256 via JWKS \\
$\square$ & Adapter \texttt{SecurityConfig} \\
$\square$ & Provisionner les templates de rôles FleetMan \\
$\square$ & Migrer les utilisateurs existants \\
$\square$ & Créer \texttt{KernelOrganizationAdapter} \\
$\square$ & Modifier \texttt{FleetService} pour créer des organisations Kernel \\
$\square$ & Créer \texttt{KernelResourceAdapter} \\
$\square$ & Modifier \texttt{VehicleService} pour enregistrer dans resource-core \\
$\square$ & Créer \texttt{KernelActorAdapter} \\
$\square$ & Créer \texttt{KernelThirdPartyAdapter} \\
$\square$ & Créer \texttt{KernelFileAdapter} \\
$\square$ & Créer \texttt{KernelSettingsAdapter} \\
$\square$ & Créer \texttt{KernelEventConsumer} \\
$\square$ & Ajouter la migration SQL \texttt{020-add-kernel-references.sql} \\
$\square$ & Configurer le Circuit Breaker \\
$\square$ & Configurer les alertes de quota \\
$\square$ & Tests d'intégration end-to-end \\
$\square$ & Documentation API mise à jour \\
\bottomrule
\caption{Checklist complète d'intégration Kernel dans FleetMan}
\end{longtable}

\end{document}
